% VI24_DETAILED_EXERCISE_SOLUTION_V1_01
\paragraph{Exercise VI.24.1.}
\begin{solution}
Let $f:X\to S$ and $g:S'\to S$.  By definition the base change is the fiber product
$X_{S'}=X\times_S S'$.  Thus it comes with projections $p_X:X_{S'}\to X$ and
$f_{S'}:X_{S'}\to S'$ fitting into
\[
\begin{tikzcd}
X_{S'} \arrow[r,"p_X"] \arrow[d,"f_{S'}"'] & X \arrow[d,"f"]\\
S' \arrow[r,"g"'] & S .
\end{tikzcd}
\]
The square is Cartesian: for every scheme $T$, giving $u:T\to X_{S'}$ is equivalent to giving
maps $a:T\to X$ and $b:T\to S'$ such that $f\circ a=g\circ b$.  The map $u$ is then uniquely
characterized by $p_Xu=a$ and $f_{S'}u=b$.  This universal property, not the drawing alone, is
what defines the base-change square.
\end{solution}

% VI24_DETAILED_EXERCISE_SOLUTION_V1_02
\paragraph{Exercise VI.24.2.}
\begin{solution}
Take the second map $S\to S$ to be $\id_S$.  A $T$-point of $X\times_S S$ is a pair
$(a:T\to X,b:T\to S)$ satisfying $f\circ a=b$.  Hence $b$ is forced once $a$ is chosen, and
there is a natural bijection
\[
\Hom(T,X\times_S S)\cong\Hom(T,X).
\]
Explicitly, the first projection $p_X:X\times_SS\to X$ is one direction.  The inverse is
$(\id_X,f):X\to X\times_SS$, since its two components are compatible over $S$.  The two composites
have the same projections as the relevant identity maps, so uniqueness in the fiber-product
universal property makes them identities.  Therefore $X\times_SS\cong X$ canonically.
\end{solution}

% VI24_DETAILED_EXERCISE_SOLUTION_V1_03
\paragraph{Exercise VI.24.3.}
\begin{solution}
Define
\[
\Phi:k[x]\otimes_k K\longrightarrow K[x],\qquad f(x)\otimes a\longmapsto a f(x),
\]
where the coefficients of $f$ are viewed in $K$.  This map is $K$-linear and multiplicative.
Conversely, the $k$-algebra map $k[x]\to K[x]$ together with the scalar map $K\to K[x]$ gives,
by the tensor-product universal property, a map $K[x]\to k[x]\otimes_kK$ sending
$x\mapsto x\otimes1$ and $a\mapsto1\otimes a$.  These maps are inverse on the generators
$K$ and $x$.  Hence
\[
k[x]\otimes_kK\cong K[x].
\]
Geometrically, $\mathbb A^1_k$ becomes $\mathbb A^1_K$ after extending the base field.
\end{solution}

% VI24_DETAILED_EXERCISE_SOLUTION_V1_04
\paragraph{Exercise VI.24.4.}
\begin{solution}
Start from the exact presentation
\[
k[x]\xrightarrow{\;\cdot(x^2+1)\;}k[x]\longrightarrow k[x]/(x^2+1)\longrightarrow0.
\]
Tensoring over the field $k$ with $K$ is exact, and Exercise VI.24.3 identifies
$k[x]\otimes_kK$ with $K[x]$.  Therefore the cokernel is
\[
(k[x]/(x^2+1))\otimes_kK\cong K[x]/(x^2+1).
\]
The polynomial is the same formal equation, but its factorization can change with $K$.  For
example, when $k=\mathbb R$ and $K=\mathbb C$,
$x^2+1=(x-i)(x+i)$ and the Chinese remainder theorem yields $\mathbb C\times\mathbb C$.
Thus scalar extension can reveal components invisible over the smaller field.
\end{solution}

% VI24_DETAILED_EXERCISE_SOLUTION_V1_05
\paragraph{Exercise VI.24.5.}
\begin{solution}
The quotient map $B\twoheadrightarrow B/I$ gives a right-exact sequence of $A$-modules
\[
I\longrightarrow B\longrightarrow B/I\longrightarrow0.
\]
Tensor with the $A$-algebra $A'$.  Right exactness gives
\[
I\otimes_AA'\longrightarrow B\otimes_AA'\longrightarrow(B/I)\otimes_AA'\longrightarrow0.
\]
Hence $(B/I)\otimes_AA'$ is the quotient of $B\otimes_AA'$ by the image of
$I\otimes_AA'$.  This image is precisely the ideal generated by the elements $i\otimes1$ with
$i\in I$, usually denoted $I(B\otimes_AA')$.  Therefore
\[
(B/I)\otimes_AA'\cong(B\otimes_AA')/I(B\otimes_AA').
\]
This is the ring-theoretic reason a closed subscheme remains closed after base change.
\end{solution}

% VI24_DETAILED_EXERCISE_SOLUTION_V1_06
\paragraph{Exercise VI.24.6.}
\begin{solution}
Let $b\in B$.  Consider the canonical map
\[
B\otimes_AA'\longrightarrow B_b\otimes_AA'.
\]
The element $b\otimes1$ becomes invertible on the right, so the universal property of localization
produces a map
\[
(B\otimes_AA')_{b\otimes1}\longrightarrow B_b\otimes_AA'.
\]
In the other direction, the maps $B\to(B\otimes_AA')_{b\otimes1}$ and
$A'\to(B\otimes_AA')_{b\otimes1}$ are $A$-compatible, and the image of $b$ is invertible.  Thus
$B\to\cdots$ factors through $B_b$, giving the inverse map.  Consequently
\[
B_b\otimes_AA'\cong(B\otimes_AA')_{b\otimes1}.
\]
Taking spectra says that the pullback of $D(b)$ is $D(b\otimes1)$.
\end{solution}

% VI24_DETAILED_EXERCISE_SOLUTION_V1_07
\paragraph{Exercise VI.24.7.}
\begin{solution}
Write $\mathbb A^1_k=\Spec k[x]$ and $D(x)=\Spec k[x,x^{-1}]$.  After the field extension
$k\to K$,
\[
k[x,x^{-1}]\otimes_kK\cong K[x,x^{-1}]
\]
by localization compatibility.  Also $k[x]\otimes_kK\cong K[x]$.  Therefore the base-changed
open immersion is
\[
D(x)_K=\Spec K[x,x^{-1}]\hookrightarrow\Spec K[x]=\mathbb A^1_K,
\]
which is exactly the principal open $D(x)$ over $K$.  Thus the same nonvanishing condition survives
scalar extension.
\end{solution}

% VI24_DETAILED_EXERCISE_SOLUTION_V1_08
\paragraph{Exercise VI.24.8.}
\begin{solution}
The closed subscheme $V(x^2)\subseteq\mathbb A^1_k$ is
$\Spec(k[x]/(x^2))$.  Tensoring with $K$ gives
\[
(k[x]/(x^2))\otimes_kK\cong K[x]/(x^2).
\]
Hence its base change is
\[
V(x^2)_K=\Spec K[x]/(x^2)\hookrightarrow\mathbb A^1_K.
\]
The underlying topological support is still the origin, while the class of $x$ remains a nonzero
square-zero element.  So base change preserves this nilpotent thickening; it does not replace it by
the reduced point $V(x)$.
\end{solution}

% VI24_DETAILED_EXERCISE_SOLUTION_V1_09
\paragraph{Exercise VI.24.9.}
\begin{solution}
Let $u:X\to Y$ be an isomorphism over $S$ with inverse $v:Y\to X$.  Base change is functorial,
so
\[
(v\circ u)_{S'}=v_{S'}\circ u_{S'},\qquad
(u\circ v)_{S'}=u_{S'}\circ v_{S'}.
\]
Since $v\circ u=\id_X$ and $u\circ v=\id_Y$, and identity morphisms base-change to identities,
we obtain
\[
v_{S'}u_{S'}=\id_{X_{S'}},\qquad
u_{S'}v_{S'}=\id_{Y_{S'}}.
\]
Therefore $u_{S'}$ is an isomorphism with inverse $v_{S'}$.  This argument is purely formal and
requires no affine coordinates.
\end{solution}

% VI24_DETAILED_EXERCISE_SOLUTION_V1_10
\paragraph{Exercise VI.24.10.}
\begin{solution}
Let $A\to A'\to A''$ and let $B$ be an $A$-algebra.  First base-change to $A'$ and then to $A''$:
\[
(B\otimes_AA')\otimes_{A'}A''.
\]
Associativity of tensor products gives a canonical map
\[
(B\otimes_AA')\otimes_{A'}A''\longrightarrow B\otimes_AA'',
\quad (b\otimes a')\otimes a''\longmapsto b\otimes a'a''.
\]
Its inverse sends $b\otimes a''$ to $(b\otimes1)\otimes a''$.  Hence
\[
(B\otimes_AA')\otimes_{A'}A''\cong B\otimes_AA''.
\]
Taking spectra yields $(X_{S'})_{S''}\cong X_{S''}$ in the affine case, and the general statement
follows from the universal property.
\end{solution}

% VI24_DETAILED_EXERCISE_SOLUTION_V1_11
\paragraph{Exercise VI.24.11.}
\begin{solution}
Identify $\mathbb R[x]/(x^2+1)$ with $\mathbb C$.  Scalar extension gives
\[
\mathbb C\otimes_{\mathbb R}\mathbb C
\cong \mathbb C[x]/(x^2+1).
\]
Over $\mathbb C$ the polynomial factors as $(x-i)(x+i)$, and the two ideals are comaximal because
$2i$ is a unit.  The Chinese remainder theorem therefore gives
\[
\mathbb C[x]/((x-i)(x+i))
\cong\mathbb C[x]/(x-i)\times\mathbb C[x]/(x+i)
\cong\mathbb C\times\mathbb C.
\]
Thus a one-point $\mathbb R$-scheme with residue field $\mathbb C$ becomes a disjoint union of two
complex points after base change to $\mathbb C$.
\end{solution}

% VI24_DETAILED_EXERCISE_SOLUTION_V1_12
\paragraph{Exercise VI.24.12.}
\begin{solution}
Use the isomorphism from Exercise VI.24.11,
\[
\mathbb C\otimes_{\mathbb R}\mathbb C\cong\mathbb C\times\mathbb C.
\]
The elements $(1,0)$ and $(0,1)$ satisfy $e^2=e$, are neither $0$ nor $1$, and sum to $1$.
Transporting them through the isomorphism gives nontrivial idempotents in the tensor product.  More
explicitly, if $i$ denotes the imaginary unit, one may take
\[
e_+=\frac12(1\otimes1-i\otimes i),\qquad
e_-=\frac12(1\otimes1+i\otimes i).
\]
Using $(i\otimes i)^2=1\otimes1$ verifies $e_\pm^2=e_\pm$, $e_+e_-=0$, and
$e_++e_-=1$.  Algebraically these idempotents encode the two connected components of the spectrum.
\end{solution}

% VI24_DETAILED_EXERCISE_SOLUTION_V1_13
\paragraph{Exercise VI.24.13.}
\begin{solution}
Let $B=\mathbb Z[x]/(x^2-2)$.  Base change along $\mathbb Z\to\mathbb Q$ gives
\[
B\otimes_{\mathbb Z}\mathbb Q
\cong \mathbb Q[x]/(x^2-2),
\]
using the polynomial-quotient scalar-extension formula.  Since $x^2-2$ is irreducible over
$\mathbb Q$, this quotient is the field $\mathbb Q(\sqrt2)$.  Thus the generic characteristic-zero
base change is a one-point affine scheme with residue field $\mathbb Q(\sqrt2)$.
\end{solution}

% VI24_DETAILED_EXERCISE_SOLUTION_V1_14
\paragraph{Exercise VI.24.14.}
\begin{solution}
Now base-change $B=\mathbb Z[x]/(x^2-2)$ along $\mathbb Z\to\mathbb F_2$.  We obtain
\[
B\otimes_{\mathbb Z}\mathbb F_2
\cong\mathbb F_2[x]/(x^2-\overline2)
=\mathbb F_2[x]/(x^2).
\]
The class $\bar x$ is nonzero but satisfies $\bar x^2=0$, so the scalar extension is nonreduced.
This computation already shows that changing characteristic can alter scheme structure sharply;
VI/25 will reinterpret the same ring as the fiber over the prime $(2)$ and analyze such fibers
systematically.
\end{solution}

% VI24_DETAILED_EXERCISE_SOLUTION_V1_15
\paragraph{Exercise VI.24.15.}
\begin{solution}
Let $X,Y\to S$ and $S'\to S$.  For a test scheme $T\to S'$, a map
\[
T\to (X\times_SY)\times_SS'
\]
is equivalent to maps $a:T\to X$, $b:T\to Y$, and $c:T\to S'$ such that the composites of
$a,b,c$ to $S$ agree.  On the other hand, a map
\[
T\to X_{S'}\times_{S'}Y_{S'}
\]
is a pair of maps to $X_{S'}$ and $Y_{S'}$ with the same $S'$-component, which unfolds to exactly
the same triple $(a,b,c)$.  The identifications are natural in $T$, hence
\[
(X\times_SY)_{S'}\cong X_{S'}\times_{S'}Y_{S'}.
\]
\end{solution}

% VI24_DETAILED_EXERCISE_SOLUTION_V1_16
\paragraph{Exercise VI.24.16.}
\begin{solution}
Base-change the diagonal $\Delta_{X/S}:X\to X\times_SX$.  By Exercise VI.24.15,
\[
(X\times_SX)_{S'}\cong X_{S'}\times_{S'}X_{S'}.
\]
The base-changed morphism $X_{S'}\to(X\times_SX)_{S'}$ has both compositions with the two
projections equal to $\id_{X_{S'}}$, because the original diagonal has both components equal to
$\id_X$.  The diagonal $\Delta_{X_{S'}/S'}$ is uniquely characterized by that same property.
Therefore
\[
(\Delta_{X/S})_{S'}\cong\Delta_{X_{S'}/S'}.
\]
\end{solution}

% VI24_DETAILED_EXERCISE_SOLUTION_V1_17
\paragraph{Exercise VI.24.17.}
\begin{solution}
For $u:X\to Y$ over $S$, the graph $\Gamma_u:X\to X\times_SY$ is characterized by the two
components $(\id_X,u)$.  After base change, the target becomes
\[
(X\times_SY)_{S'}\cong X_{S'}\times_{S'}Y_{S'},
\]
and the two components of the base-changed graph are $\id_{X_{S'}}$ and $u_{S'}$.  These are
exactly the defining components of $\Gamma_{u_{S'}}$.  Hence
\[
(\Gamma_u)_{S'}\cong\Gamma_{u_{S'}}.
\]
Equivalently, use the fact that a graph is a pullback of a diagonal and that pullbacks compose.
\end{solution}

% VI24_DETAILED_EXERCISE_SOLUTION_V1_18
\paragraph{Exercise VI.24.18.}
\begin{solution}
Connectedness and irreducibility need not be preserved by field extension.  Take
$X=\Spec\mathbb C$ as an $\mathbb R$-scheme.  It has a single point, so it is connected and
irreducible.  After base change to $\mathbb C$,
\[
X_{\mathbb C}=\Spec(\mathbb C\otimes_{\mathbb R}\mathbb C)
\cong\Spec(\mathbb C\times\mathbb C),
\]
which is a disjoint union of two points.  It is therefore disconnected and reducible.  The
example shows that base change preserves the Cartesian universal property, not arbitrary global
topological properties.
\end{solution}

% VI24_DETAILED_EXERCISE_SOLUTION_V1_19
\paragraph{Exercise VI.24.19.}
\begin{solution}
Let $p>0$, $k=\mathbb F_p(s^p)$ and $L=\mathbb F_p(s)$.  The extension is purely inseparable of
degree $p$.  Introduce a variable $u$ for the second copy of $s$; then
\[
L\cong k[u]/(u^p-s^p).
\]
After extending scalars from $k$ to $L$,
\[
L\otimes_kL\cong L[u]/(u^p-s^p).
\]
In characteristic $p$, $u^p-s^p=(u-s)^p$, so
\[
L\otimes_kL\cong L[u]/((u-s)^p).
\]
The class of $u-s$ is nonzero and nilpotent.  Thus two reduced one-point schemes can have a
nonreduced fiber product after an inseparable field extension.
\end{solution}

% VI24_DETAILED_EXERCISE_SOLUTION_V1_20
\paragraph{Exercise VI.24.20.}
\begin{solution}
Work affine-locally on the target.  A closed immersion is represented by a surjective ring map
$B\twoheadrightarrow B/I$.  After base change along $A\to A'$, its coordinate map is
\[
B\otimes_AA'\longrightarrow(B/I)\otimes_AA'.
\]
By right exactness of tensor product, this map is again surjective, and Exercise VI.24.5 identifies
the target as
\[
(B\otimes_AA')/I(B\otimes_AA').
\]
Thus the pullback is again an affine closed subscheme.  Since being a closed immersion is local on
the target, the general base-changed morphism is a closed immersion.
\end{solution}

% VI24_DETAILED_EXERCISE_SOLUTION_V1_21
\paragraph{Exercise VI.24.21.}
\begin{solution}
Again work locally.  An open immersion is locally a principal-open inclusion
$D(b)=\Spec B_b\hookrightarrow\Spec B$.  After base change along $A\to A'$, Exercise VI.24.6 gives
\[
B_b\otimes_AA'\cong(B\otimes_AA')_{b\otimes1}.
\]
Therefore the pullback of $D(b)$ is the principal open $D(b\otimes1)$ of
$\Spec(B\otimes_AA')$.  Principal opens form a basis, so these local descriptions glue.  Hence an
arbitrary open immersion remains an open immersion after arbitrary base change.
\end{solution}

% VI24_DETAILED_EXERCISE_SOLUTION_V1_22
\paragraph{Exercise VI.24.22.}
\begin{solution}
Let $u:X\to Y$ be a monomorphism over $S$, and let $u_{S'}:X_{S'}\to Y_{S'}$ be its base change.
Suppose $a,b:T\to X_{S'}$ satisfy $u_{S'}a=u_{S'}b$.  Compose with the projection
$p_X:X_{S'}\to X$.  The equality after applying $u_{S'}$ implies
\[
u\,p_Xa=u\,p_Xb.
\]
Since $u$ is a monomorphism, $p_Xa=p_Xb$.  The second component of both $a$ and $b$ is already the
fixed structure map $T\to S'$ because they are morphisms over $S'$.  A map to the fiber product
$X\times_SS'$ is uniquely determined by these two components, so $a=b$.  Thus $u_{S'}$ is a
monomorphism.
\end{solution}

% VI24_DETAILED_EXERCISE_SOLUTION_V1_23
\paragraph{Exercise VI.24.23.}
\begin{solution}
Let $B=k[x,y]/(xy)$.  Tensor the exact quotient presentation with $K$:
\[
B\otimes_kK
\cong (k[x,y]\otimes_kK)/(xy)
\cong K[x,y]/(xy).
\]
Thus the scalar extension is the same crossing equation, now over $K$:
\[
(\Spec k[x,y]/(xy))_K\cong\Spec K[x,y]/(xy).
\]
The two coordinate axes remain the two closed components when the base field is enlarged.  The
calculation illustrates that coefficients extend while the defining ideal is extended to the new
polynomial ring.
\end{solution}

% VI24_DETAILED_EXERCISE_SOLUTION_V1_24
\paragraph{Exercise VI.24.24.}
\begin{solution}
Since $\mathbb A^2_k=\Spec k[x,y]$, affine base change gives
\[
\mathbb A^2_k\times_{\Spec k}\Spec K
=\Spec(k[x,y]\otimes_kK).
\]
As in the one-variable case,
\[
k[x,y]\otimes_kK\cong K[x,y],
\]
with $x\mapsto x\otimes1$ and $y\mapsto y\otimes1$.  Therefore
\[
(\mathbb A^2_k)_K\cong\mathbb A^2_K.
\]
More generally the same argument gives $(\mathbb A^n_k)_K\cong\mathbb A^n_K$.
\end{solution}

% VI24_DETAILED_EXERCISE_SOLUTION_V1_25
\paragraph{Exercise VI.24.25.}
\begin{solution}
For a point $s\in S$, the residue field determines a canonical morphism
\[
\Spec\kappa(s)\longrightarrow S
\]
whose image is $s$.  Base-changing $X\to S$ along this morphism gives
\[
X_s:=X\times_S\Spec\kappa(s).
\]
This is the scheme-theoretic fiber over $s$.  Formally it is nothing more than a special instance
of the VI/24 base-change construction.  VI/25 studies what this particular base change looks like:
its coordinate rings, points, components, residue fields, and geometric variants.
\end{solution}

% VI24_DETAILED_EXERCISE_SOLUTION_V1_26
\paragraph{Exercise VI.24.26.}
\begin{solution}
VI/24 develops the general operation
\[
X\longmapsto X_{S'}=X\times_SS'
\]
for an arbitrary change of base $S'\to S$, together with its functoriality and elementary
stability properties.  VI/25 specializes this operation to $S'=\Spec\kappa(s)$ (and field
extensions thereof) and studies the geometry of the resulting fibers.  Thus VI/24 is the calculus
of changing the base; VI/25 is the geometry of the special base changes selected by points.
\end{solution}
